The mean-square estimate \eqref{eq:intro-ms} gives a power saving because $A_{b^6}(D)$ is, on average over $b\in\OO$, close to a $\nu$-twisted M\"obius sum. Indeed, the identity $\chi_n(b^6)=\ind_{(b,n)=1}$ shows that averaging $A_{b^6}(D)$ over $\NK(b)\le Y$ against a smooth bump $V$ counts, for each $n$, the $b$ prime to $n$, of which there are about $Y\varphi(n)/\NK(n)$. Poisson summation in $b$ makes this precise:
\begin{equation}\label{eq:intro-average}
 \sum_{b\in\OO}V\Bigl(\frac{\NK(b)}{Y}\Bigr)A_{b^6}(D)=Y\widehat V(0)A^\varphi(D)+O_{W,V,\varepsilon}(D^{1+\varepsilon}),
\end{equation}
where $\widehat V(0)\ne0$ and
\[
 A^\varphi(D)=\sum_{\mathfrak n}\mu(\mathfrak n)\nu(\mathfrak n)\frac{\varphi(\mathfrak n)}{\NK(\mathfrak n)}W(\NK(\mathfrak n)/D),
 \qquad
 \frac{\varphi(\mathfrak n)}{\NK(\mathfrak n)}=\prod_{\mathfrak p\mid\mathfrak n}\Bigl(1-\frac1{\NK(\mathfrak p)}\Bigr).
\]
Taking $Y=H^{1/6}$, so that the rows $b^6$ have norm at most $H$, and applying Cauchy--Schwarz and \eqref{eq:intro-ms} to the left side gives
\begin{equation}\label{eq:intro-extraction}
|A^\varphi(D)|^2
\ll D^{1+\varepsilon}H^{5/6}+D^{2+2\varepsilon}H^{-1/3}.
\end{equation}
With $H=D^{1+\vartheta}$, this gives
$A^\varphi(D)\ll D^{11/12+5\vartheta/12+\varepsilon}=D^{23/24+\varepsilon}$, the same exponent that the rows $p^6$ give in \cite{OpenAI26b}. The weight $\varphi(\mathfrak n)/\NK(\mathfrak n)$ changes the Dirichlet series of the M\"obius sum only by an Euler product that converges absolutely in $\Rea s>0$, so this saving serves as well as \eqref{eq:power-savings}.
Therefore it remains to establish \eqref{eq:intro-ms}. A smaller $\vartheta$ would give a better exponent; we take $\vartheta=1/10$ because the formal proof of \eqref{eq:intro-ms} needs $\vartheta>1/20$ (Remark~\ref{rem:formal-descent}).

